\chapter{An Executable Reference Interpreter}
\label{ch:pebble-interpreter}

A representation specification becomes easier to challenge when it can run.
The interpreter is our first executable consumer of every operation and terminator.
It is deliberately direct: verification first, one register slot per value, explicit phi capture, and explicit dispatch.
It is not a performance model of the final machine.

\section{Execution without hidden control state}
\label{sec:pebble-interpreter-loop}
\code{pir\_run} accepts a function, a byte input, and a block-entry budget.
It rejects out-of-domain inputs and malformed functions.
Registers begin at zero as a host initialization measure, but valid execution must not depend on those defaults.
Dominance and phi-edge checks ensure that each used value comes from a definition on the executed path.

Algorithm~\ref{alg:pebble-run} follows the block semantics from Chapter~\ref{ch:pebble-semantics}.
Each phi prefix reads from the register state established before that block entry.
Only after all incoming values have been captured are the phi result registers updated.

\begin{algorithm}[htbp]
\caption{Interpreting a verified function with an explicit block budget.}
\label{alg:pebble-run}
\begin{pseudocode}
\Require function $F$, byte input $x$, nonnegative budget $K$
\Ensure return value, invalid-program result, or budget exhaustion
\State verify $F$; reject on failure
\State $B\gets0$; $P\gets$ no predecessor; $steps\gets0$
\While{$steps<K$}
\State increment $steps$
\State capture phi inputs of $B$ selected by $P$ into scratch storage
\State publish phi results and evaluate ordinary instructions of $B$
\If{the terminator returns $v$} \State \Return the current value of $v$ \EndIf
\State $P\gets B$; choose the successor of $B$ from its terminator
\EndWhile
\State \Return budget exhaustion
\end{pseudocode}
\end{algorithm}

\section{Three independent traces}
\label{sec:pebble-interpreter-traces}
For the running expression at input 9, values 1 through 7 are 9, 7, 16, 3, 5, 8, and 128.
The function returns after one block entry.
After optimization, values 4 and 5 are dead and value 6 is directly 8; the observable answer and block count remain unchanged.

For the diamond at input 12, entry chooses the right arm.
That arm computes 24, and the join phi selects its operand labeled by predecessor 2.
The return is 24 after three block entries.
The unused left-arm slots' host initialization must not influence the selected result.

For the summation loop at input 3, the dynamic block sequence is
$0,1,2,1,2,1,2,1,3$.
There are nine block entries, three body executions, and result 3.
A budget of eight exhausts just before entering the exit block; a budget of nine permits the return.
That off-by-one case is important because a budget measures interpreter transitions, not source-loop iterations.

\section{Reference does not mean infallible}
\label{sec:pebble-oracles}
The IR interpreter and optimizer share the representation, so their agreement alone can conceal a shared misunderstanding.
The generated-expression tests additionally evaluate a separate source tree using its own recursion.
The loop tests compare with direct C computations of the mathematical recurrence, reduced to bytes.
The lowered VM provides a different execution mechanism, including explicit edge-copy sequences.
Agreement among these paths is stronger evidence than rerunning one interpreter under two names.

Still, finite agreement is not a general proof.
The tests enumerate all 256 input bytes for each chosen program, but they do not enumerate all possible CFGs or expressions.
The local semantic arguments remain necessary, especially when an extension introduces memory, traps, or nondeterministic observations.

\section{Cost and interface boundaries}
\label{sec:pebble-run-cost}
With $N$ allocated slots and $K$ executed block entries, the reference interpreter uses $O(KN)$ table-scan work after verification.
Register and scratch arrays use $O(N)$ storage.
Each ordinary operation takes constant time for this fixed byte domain.
The interpreter's fuel bounds its own dispatch loop because each block scan is bounded and its operations terminate.
That reasoning would need revision if an instruction could call an unbounded external function.

The VM counts individual instructions instead of block entries.
Comparisons between the two interpreters must use budgets known to be sufficient or distinguish exhaustion from semantic disagreement.
Equal numeric budgets are not equal work and do not define a cross-interpreter equivalence relation.

\section{Exercises}
\label{sec:pebble-twelve-exercises}
\begin{enumerate}
\item \label{ex:pebble-loop-budget} Derive the summation loop's block-entry count as a function of input $x$, including entry, the final header visit, and exit.
\item \label{ex:pebble-zero-registers} Why does initializing every register to zero not make a missing definition harmless? Construct a result that would be silently wrong without verification.
\end{enumerate}

\section{Further reading}
\label{sec:pebble-twelve-reading}
Appel's executable approach to language implementation is useful when designing a small semantic oracle~\cite{appel-1998-modern}.
Safonov helps distinguish validation evidence from a complete correctness claim~\cite{safonov-2010-trustworthy}.
